%!TEX program = xelatex
\documentclass[a4paper,12pt]{article}
\usepackage{fontspec}\defaultfontfeatures{Ligatures=TeX}
\usepackage[a4paper,vmargin={3cm,3.5cm},hmargin={3cm,3cm}]{geometry}
%-----------------------------------------------------------------------------%
\usepackage{settings}
%-----------------------------------------------------------------------------%
%%% Title %%%
\title{G6411 Recitation Note: Defining Terms}
\author{Jesse Chieh Chen\\\href{mailto:cc5229@columbia.edu}{cc5229@columbia.edu}}
\date{\today}
%-----------------------------------------------------------------------------%

\begin{document}

\maketitle

\noindent
Let $X$ be a real-valued random variable.

\begin{definition}[\glsentryshort{cdf} and \glsentryshort{pdf}]\label{def-cdf-pdf}
	The \gls{cdf} of $X$ is defined as $F_X(x) \coloneqq \Pr\{X \leq x\}$ for all $x \in \reals$.
	A \gls{pdf} of $X$ is a nonnegative measurable function $f_X$ satisfying
	$F_X(x)=\int_{-\infty}^{x}f_X(t)\diff t$ for all $x\in\reals$.
	Such a density exists if and only if $F_X$ is absolutely continuous.
	In that case, $f_X(x)=F_X'(x)$ for almost every $x$.
\end{definition}

\begin{remark}
	The distribution of $X$ is characterized by its \gls{cdf},
	so the \gls{cdf} exists for any random variable,
	whereas the \gls{pdf} may not exist for all random variables.
	For example, the \gls{cdf} of a discrete random variable has jumps and cannot be represented as an integral of a density.
\end{remark}

\begin{definition}[Moment]\label{def-moment}
	The $k$-th moment of a random variable $X$ is defined as $\E(X^k)$, provided that the expectation exists.
\end{definition}

\begin{remark}
	The mean is the first moment,
	and the variance is the second \emph{central} moment,
	defined as $\var(X) = \E[(X - \E X)^2]$.
	The second moment of $X$ is $\E X^2$.
\end{remark}

\begin{definition}[Sample]\label{def-sample}
	A sample is a collection of random variables,
	often denoted as $\{X_i\}_{i=1}^{n}$ where $n$ is the sample size.
\end{definition}

\begin{definition}[Statistic]\label{def-statistic}
	A statistic is a function of a sample.
\end{definition}

\begin{definition}[Estimator]\label{def-estimator}
	An estimator is a statistic used to guess, compute, or estimate a parameter.
\end{definition}

% Three estimators
% - analog: MoM
% - loss function: square loss
% - fully parametric
%     - likelihood based methods

\begin{remark}
	An \emph{estimate} is the realized value of an estimator given a realized sample.
	An \emph{estimand} is the parameter that an estimator aims to estimate.
\end{remark}

\begin{definition}[Parameter]\label{def-parameter}
	A parameter is a target feature of the underlying model.
\end{definition}

\begin{remark}[What is a parameter really?]
	There are different ways to define a parameter.
	In class, we defined a parameter as a function of the population distribution,
	which is a definition commonly used in statistics.
	Under this definition, identification is built in:
	once the population distribution is given, the parameter is uniquely determined.
	% This does not mean that we know its value from a finite sample.
	In econometrics, we also work with \emph{economic} models,
	whose parameters may describe preferences, technologies, or behavioral responses.
	Different values of these parameters can sometimes generate the same distribution of observable variables,
	so identification becomes a separate question.
\end{remark}

\begin{definition}[\glsentrylong{iid}]\label{def-iid}
	A sample $\{X_i\}_{i=1}^{n}$ is said to be \gls{iid} if the $X_i$ are independent and identically distributed,
	denoted as $\{X_i\}_{i=1}^n\iidto F$ for some distribution $F$.
\end{definition}

\begin{remark}
	Here is an example of a sample that is identically distributed but \emph{not} independent:
	Let $\{Z_i\}_{i=0}^n\iidto\normaldistribution(0,1)$, with $n\geq2$.
	Define $X_i = (Z_0 + Z_i)/\sqrt{2}$ for all $i=1,...,n$.
	Then each $X_i\sim\normaldistribution(0,1)$,
	but $\cov(X_i,X_j)=1/2$ for $i\ne j$ because they share the term $Z_0$.
	Thus, the $X_i$ are identically distributed but not independent.
\end{remark}

\begin{definition}[Identification]\label{def-identification}
	A parameter is said to be identified if it can be written as a function of the population distribution,
	i.e., if it can be uniquely determined from the population distribution of the observable variables.
\end{definition}

\begin{remark}[Alternative Formulation]
	You will often see identification formulated in the following way:
	Let $\Theta$ denote the parameter space and let $P_\theta$ denote the population distribution of the observable variables under parameter $\theta\in\Theta$.
	Then the parameter $\theta$ is said to be identified if $P_\theta=P_{\theta'}$ implies $\theta=\theta'$ for all $\theta,\theta'\in\Theta$.
\end{remark}

\begin{example}[Schooling and Wages]\label{eg-identification-schooling}
	Suppose that log wages $Y$ are determined by years of schooling $S$ and an unobserved factor $U$:
	\begin{align}
		Y=\alpha+\beta S+U,
		\qquad \E(U)=0.
	\end{align}
	Here $\beta$ describes the effect of an additional year of schooling holding $U$ fixed.
	We observe $(Y,S)$, but not $U$, which may include unobserved ability.
	Assume finite first moments and place no further restrictions on the relationship between $S$ and $U$.
	For any candidate value $b$, define
	\begin{align*}
		a_b&=\E(Y)-b\E(S),\\
		U_b&=Y-a_b-bS.
	\end{align*}
	Then $Y=a_b+bS+U_b$ and $\E(U_b)=0$.
	Thus, every value $b$ is compatible with exactly the same joint distribution of the observables $(Y,S)$.
	The schooling effect $\beta$ is not identified under these assumptions:
	Even knowing the entire population distribution would not distinguish it from the contribution of the unobserved factor.
	This demonstrates that only requiring $\E(U)=0$ is not enough to identify $\beta$,
	since an arbitrary relationship between $S$ and $U$ is allowed.
\end{example}

\vfill
\section*{Acronyms}
\printglossaries

\end{document}
